\section{Conclusion}
\label{a:sec:conclusion}

The constant $A(p)$ organises the quasi-stationary mean of subcritical binary
Galton--Watson branching. It is well defined, computable, and linked to the basin
Koenigs value $A(p)=2\lim_n r^{-n}f_r^{\circ n}(\tfrac12)$ at multiplier $r=2p$.
For every $r\in(0,1)$ that conjugacy is hypertranscendental in its own variable
(\cref{a:thm:mainht}) --- a theorem with no exceptional parameters in the
attracting window, since differentially algebraic Schr\"oder functions exist only
over repelling fixed points --- and this explains the failure of every elementary
closed-form search of \cref{a:sec:GA}, while leaving $A(p)$ fully usable for
analysis and numerics through the scheme of \cref{a:sec:practical}.

What the theorem does not decide has been flagged with equal care. Transcendence
--- even irrationality --- of individual values $A(p_0)$ remains open, because the
1993 Becker--Bergweiler values theorem \cite{Becker1993} is confined to
equal-degree B\"ottcher pairs and does not transfer to the Schr\"oder setting. The
elementarity and differential algebraicity of the map $p\mapsto A(p)$, as opposed
to the Koenigs function $\psi_r$ in its argument, are likewise open; the former
question appears nowhere in the literature, and the natural machinery of
parametrised differential Galois theory \cite{HardouinSinger2008} does not apply
as it stands. Eleven rational values nonetheless resist all inverse-symbolic
identification at over $200$-digit rigour, which is what the evidence of
\cref{a:sec:pslq} amounts to --- and no more.

The exceptional elementary linearisations at $r=2$ and $r=4$, both repelling, lie
outside the branching window; their derivations, and their place inside the
Becker--Bergweiler classification, are collected in \cref{a:app:integrable}, and
the exploratory formulae of the closed-form search in \cref{a:app:catalogue}. The
open problems proper to this chapter are therefore three, corresponding to the
readings of \cref{a:sec:scope}: the arithmetic nature of the values $A(p_0)$ at
rational $p_0$; the elementarity and differential algebraicity of the parameter
map $p\mapsto A(p)$; and its expressibility, if any, in terms of standard special
functions.

\FloatBarrier
